\subsection{Complete rectangular cancellation by sparse column peeling}
\label{marked-cut:column-peeling}

Reports 95 and 97 independently sharpen the square-only discussion above.
Report 95 derives the full-row criterion from report 56's permanent-one
argument; report 97 supplies a direct stopping-set persistence proof and a
sparse native implementation. We integrate the latter queue and independent
timestamp checker, extend them to the whole-cut boundary kernel, and retain
the old paths for differential measurement and certificate compatibility.
The stopping-set and iterative peeling framework is classical in coding
theory; its maximal residual is described by Di et al.\cite{dptru2002}.
The application here uses unsigned integral incidences and the geometric
positive-pivot rule, not parity-check arithmetic.

This corrects our earlier description of the capped planner as incomplete
for rectangular full-row cancellation. It was already complete for that
precise problem because it prioritizes zero-fill pivots. The improvement
removes repeated matrix scans and makes the obstruction explicit. Neither
algorithm is complete for recognition of all geometric handlebodies.

\paragraph{Definition.}
For a nonnegative integral matrix $N$, a nonempty row set $R$ is a
\emph{stopping set} if every column total
$S_h(R)=\sum_{i\in R}n_{ih}$ is zero or at least two.
A \emph{unit column} has total one on all currently live rows.
Its unique incident row has entry one; cancelling this pair creates no fill
in the other rows.


\begin{proposition}[Stopping-set persistence]
\label{marked-cut:stopping-persistence}
A nonempty stopping set prevents every full positive-pivot cancellation,
including sequences that use nonzero fill.
\end{proposition}
\begin{proof}
Pivot at $(r,g)$ with $n_{rg}=1$. If $r\notin R$, then for surviving $h$,
\[
 S'_h(R)=S_h(R)+S_g(R)n_{rh},
\]
which is zero or at least two. If $r\in R$, the stopping condition implies
$S_g(R)\ge2$, so $R'=R\setminus\{r\}$ is nonempty. Then
\[
 S'_h(R')=S_h(R)+(S_g(R)-2)n_{rh}.
\]
If $S_h(R)=0$, also $n_{rh}=0$; otherwise the right side remains at least
two. A nonempty stopping set therefore survives every pivot prefix and
prevents deletion of all rows. Saturation at two preserves this argument's
zero and singleton predicates.
\end{proof}

\begin{proposition}[Complete full-row decision]
\label{marked-cut:rectangular-complete}
For every finite rectangular nonnegative integral matrix, a full positive
cancellation exists if and only if maximal unit-column peeling deletes all
rows. If peeling stops, its residual rows form the unique largest source
stopping set, independent of queue order.
\end{proposition}
\begin{proof}
Every peel is a valid positive pivot with zero fill. If peeling stops with
rows $R$, each remaining column has total zero or at least two on $R$.
A deleted column had only its deleted row present, hence is zero on $R$.
No fill occurred, so $R$ is a stopping set in the original source matrix.
The preceding proposition excludes any other full positive schedule.

The union of stopping sets is a stopping set: any column meeting the union
already meets one constituent with total at least two. No unit-column peel
can remove the first row of any original stopping set. Every such set lies
in the final residual, which is itself a stopping set. This proves maximality
and uniqueness.
\end{proof}

A zero row is an immediate obstruction. Conversely, the all-one $2\times2$
matrix is obstructed although an arbitrary raw pivot removes one row and
leaves the entry two. The canonical residual is therefore not a maximum
partial-cancellation answer. Nor does a stopping set prove knottedness or
non-handlebody topology: a changed forest, new disk system, two/three
cancellation, or geometric intersection reduction changes the problem.

\paragraph{Why the old producer was complete.}
The retained capped planner minimizes
$(\sum_hn_{rh}-1)(\sum_in_{ig}-1)$ over singleton pivots. A zero score means
a unit row or unit column, and such a pivot has no fill. Whenever full
completion is possible, the theorem supplies a unit column and thus a
zero-score choice. A column peel preserves full cancellability because any
stopping set in the remainder was already a source stopping set. For a unit
row peel, a putative residual stopping set $R$ either has zero or at least
two crossings at the deleted column, making $R$ a source obstruction, or
has one crossing there, making $R\cup\{r\}$ a source obstruction. Thus
neither zero-fill choice can destroy a possible full completion. Induction
proves the old minimum-score producer's completeness. This does not assert
that every arbitrary singleton choice has the same property.

\paragraph{Sparse production and source timestamps.}
The queue stores, for each appearing column, its number of incident live
rows and the xor of their labels. Degree one and a unit entry identify an
eligible pivot. Deleting its row visits that row's immutable source columns
once. A column reaches degree one at most once; stale queue entries are
ignored. No incidence matrix is updated. For $I$ source nonzeros and $m,n$
rows and columns, the indexed algorithm needs $O(I+m+n)$ structural work
and storage. Reading the input integers and their label bit lengths is
charged separately; Python dictionaries are not an adversarial constant-time
hashing theorem.

The checker assigns a deletion time $\tau(r)$ to every scheduled row. For
pivot $(r,g)$ at time $t$, it checks $n_{rg}=1$ and requires
$\tau(i)<t$ for every other source row $i$ meeting $g$. Unlisted rows have
infinite time. Distinct pivot rows and columns, these incidence inequalities,
and coverage of all rows prove a complete no-fill schedule in one source
scan. It imports no producer and does not simulate its queue. A separate
checker accepts any nonempty stopping subset by verifying its capped source
column totals; canonical discovery is not trusted by the obstruction checker.

The native core keeps its earlier row-leaf shortcut, then uses column peeling.
Successful fallback proofs use \path{normal-core-handlebody-v2}; v1 leaf and
capped proofs remain replayable. The whole-cut v5 wrapper now emits
\path{bundle-boundary-kernel-v2} with source-timestamp replay. Earlier kernel
v1 proofs retain their dense positive-update replay. Omitted bundle and
connecting columns still have only zero or at least two entries, so the
projected full-row theorem applies without expanding them.

These are linear bounds for the cancellation and timestamp steps. The
retained row heap adds a logarithmic term, and the existing independent
forest and component bookkeeping can still take quadratic work. Source
orbit queries, binary arithmetic, source-path reconstruction and passive
v4 curve transport retain their separate costs. We do not describe the
entire normal-cut operator as linear.

\paragraph{An early source-core obstruction.}
Adding rows cannot remove a stopping subset of original rows. Adding columns
that are zero on those original rows cannot remove it either. In every
supported bundle presentation, the original core rows are unchanged after
the same core forest contraction; their new bundle and connecting columns
are zero. Hence a failed original-core full-row test permits an inconclusive
return before ordered boundary discovery. The current wrapper enables this
filter by default and reuses a successful core proof when disk absorption
adds no rows. Setting \path{precheck_core=False} retains the full query path
for differential checks. Three-handles are excluded before this argument.
This is a refusal of the fixed raw method, not a topological negative verdict.

\paragraph{A changed schedule can change passive witnesses.}
Although complete core cancellation has the same success set, a new schedule
can leave different cocore disks. A formerly private annulus may lose its
private column in that disk system. The complete joint boundary kernel then
provides another whole-cut certificate. The source audit must therefore compare
whole-cut censuses and replay all old proofs, rather than demand unchanged
certificate bytes or a fixed outer schema for every input.

\paragraph{Source compatibility and independent oracle.}
All 54 focused tests pass in 80.515 seconds. They include exhaustive comparison
of column peeling with every legal saturated-pivot schedule on 31,096 small
matrices, 350 seeded all-subset/minor checks, a 4,000-row incidence-work check,
hostile records, all 106 source cores, source-only boundary-query suppression,
legacy dense replay, and complete supplied-cut certificate replay.

The fresh 106-case audit retains 85 complete supplied-cut certificates and
21 inconclusive outcomes. All 85 replay independently; the 78 cases below
the expansion cap agree with literal Regina cuts, and seven large-coordinate
cases are excluded from literal expansion. The refusal split remains eight
raw stalls, three square raw obstructions, and ten cases with core three-handles.

The retained baseline is commit \texttt{5578c6e6629a5c63fc5b40ea14d1f49c9d5a8c8c}.
A separate old/new/precheck-disabled audit reproduces all 85 historical
certificate hashes, replays all 85 old certificates with both checkers, and
preserves every outcome and census. The enabled and disabled precheck produce
identical new positive certificates. Six failed source cores skip ordered
boundary discovery, and their stopping subsets are independently checked.
Only 44 positive certificates remain byte-identical: changed valid schedules
are expected. The pilot Fibonacci 7 and 9 single-surface cases move from outer
v4 to v5 while retaining their independently checked censuses.

The full-suite result and isolated performance measurements are recorded
below when complete. Reproducible drivers and raw evidence are retained as
\path{data/column-peeling-*}; the six unmodified old modules and their loader
are in \path{data/column-baseline/}.

